\section{Introduction}
\label{sec:intro}

Large language model (LLM) agents have demonstrated remarkable capabilities in interactive environments, from household task completion~\citep{shridhar2021alfworld} to web navigation~\citep{yao2022webshop} and scientific experimentation~\citep{wang2022scienceworld}. The ReAct framework~\citep{yao2023react}, which interleaves reasoning and acting, has become a standard paradigm for such agents. However, a persistent limitation remains: when an agent encounters a failure---such as attempting to take an object from a closed container---it lacks mechanisms to leverage similar past experiences, often repeating the same mistake across episodes.

Recent work has introduced experiential memory to address this limitation. Reflexion~\citep{shinn2023reflexion} generates free-text reflections after failed episodes and injects them at the start of subsequent trials. ExpeL~\citep{zhao2024expel} extracts cross-task insights from trajectory pools. These approaches share a common design choice: memory is \emph{injected at episode boundaries}---either at the beginning of a new trial or when transferring between tasks. While effective, this design means that an agent must complete (or abandon) an entire episode before benefiting from experiential memory, even when failure signals are available mid-execution.

We identify this as a fundamental limitation in the \emph{retrieval timing} of existing experiential memory systems. Drawing on the taxonomy proposed by~\citet{liu2025memory}, which categorizes agent memory along forms, functions, and dynamics, we observe that prior work predominantly uses \emph{passive} retrieval timing---injecting memories at fixed points regardless of the agent's current state. In contrast, we propose \emph{active, failure-triggered} retrieval: the agent monitors its own execution for failure signals and retrieves relevant memories only when needed.

We present \method{}, a failure-triggered in-loop memory system for ReAct agents. \method{} operates through three mechanisms: (1)~a failure detector that identifies both \emph{explicit} failures (environment error signals) and \emph{implicit} failures (actions that execute but fail to make progress) during execution; (2)~a hybrid retrieval system using BM25 and sentence embeddings with Reciprocal Rank Fusion (RRF) to find relevant past failure-recovery experiences; and (3)~a post-episode extractor that distills failure-recovery pairs from trajectories into structured memory entries.

We evaluate \method{} on five diverse benchmarks: ALFWorld~\citep{shridhar2021alfworld} (household tasks, 134 environments), WebShop~\citep{yao2022webshop} (e-commerce, 100 sessions), ScienceWorld~\citep{wang2022scienceworld} (formal Step~4 no-gate protocol: 111 test episodes across 12 core task types), HotPotQA~\citep{yang2018hotpotqa} (multi-hop QA, 500 questions), and ToolBench~\citep{qin2023toolllm} (API tool use, 765 queries). We study the setting of \emph{online adaptation on recurring task instances}, where an agent encounters the same environments across epochs and can accumulate experience---motivated by deployment scenarios such as household robots interacting with the same kitchen, customer service agents handling recurring query types, and lab assistants performing the same protocols. \method{} significantly improves over ReAct on ALFWorld ($p{=}0.033$). On the formal ScienceWorld Step~4 protocol, ExpeL achieves the highest success rate (49.55\%, 57.78 average raw score), while online \method{} reaches 41.44\% success rate and the highest raw score among non-ExpeL memory settings (64.25), compared with 37.84\% and 47.21 for the ReAct baseline and 36.94\% and 49.39 for corrected Reflexion. This reinforces a complementary rather than universally superior operating regime: in-loop memory improves online progress, while episode-level insight extraction can be stronger when a full collection epoch is available.

Our contributions are:
\begin{itemize}
    \item We propose failure-triggered in-loop memory retrieval for ReAct agents, enabling mid-episode course correction based on past experiences. On the formal ScienceWorld Step~4 protocol, online \method{} reaches 41.44\% success rate and 64.25 average raw score, improving over offline memory (39.64\%, 48.83), the ReAct baseline (37.84\%, 47.21), and corrected Reflexion (36.94\%, 49.39), while ExpeL reaches the highest success rate after an insight-collection epoch (49.55\%, 57.78).
    \item We introduce a structured failure-recovery memory format---storing (failed action, observation, corrective action) triples rather than abstract reflections or success trajectories. Ablation studies show this format outperforms Reflexion-style reflections by +3.7\% ($p{=}0.056$, borderline) on ALFWorld; the advantage over success trajectories is directional but not significant ($p{=}0.428$). We present these as promising design choices rather than definitively proven advantages.
    \item As implementation choices, we combine BM25 and sentence-embedding retrieval via RRF fusion with per-environment memory isolation. On ALFWorld, the hybrid approach matches BM25-only; broader evaluation of retrieval methods across benchmarks is left to future work.
    \item We evaluate on five benchmarks with strong baselines. \method{} achieves the highest success rate on ALFWorld and the highest ScienceWorld Step~4 raw score among non-ExpeL memory settings, while ExpeL achieves the highest success rate on ScienceWorld, WebShop, HotPotQA, and ToolBench. Cross-environment memory transfer experiments on three benchmarks reveal that task homogeneity is the key predictor: on homogeneous tasks (WebShop), shared memory achieves Epoch~2 performance in a single epoch; on heterogeneous API tasks (ToolBench), it causes negative transfer ($-$19.9\%), indicating that sharing should be enabled selectively.
\end{itemize}
